\documentclass[12pt,a4paper]{article}

\usepackage{cmap}
\usepackage[T2A]{fontenc}
\usepackage[utf8]{inputenc}
\usepackage[english,russian]{babel}
\usepackage[a4paper,left=2.5cm,right=2cm,top=2cm,bottom=2cm]{geometry}
\usepackage{booktabs}
\usepackage{array}
\usepackage{enumitem}
\usepackage[table]{xcolor}
\usepackage{hyperref}

\definecolor{candy}{RGB}{217,79,112}
\definecolor{choco}{RGB}{74,44,23}

\hypersetup{
  colorlinks=true,
  linkcolor=candy,
  urlcolor=candy,
  pdftitle={Конфетная фабрика — правила игры},
  pdfauthor={Конфетная фабрика}
}

\setlength{\parindent}{0pt}
\setlength{\parskip}{6pt}

\newcommand{\term}[1]{«\textbf{#1}»}

\title{\textbf{\textcolor{choco}{Конфетная фабрика}}\\[4pt]
\large правила игры-соревнования}
\date{}

\begin{document}
\maketitle
\vspace{-2.5em}

\section*{Легенда}

Идёт 2052 год. После Великого Сахарного Дефолта мировые валюты обесценились,
и человечество вернулось к единственному активу, чья ценность никогда не
вызывала сомнений, --- к \emph{конфете}. Центробанки хранят золотовалютные
резервы в карамели, ипотеку выдают под залог ирисок, а курс шоколадки
к зефиру объявляют по радио дважды в день.

Вы --- совет директоров конфетной фабрики. Ваши конвейеры круглосуточно
выпускают конфеты, но того же самого добиваются и конкуренты, поэтому
почивать на карамельных лаврах некогда. Единственный путь к господству
на рынке --- \emph{инновационные рецепты}: их разрабатывают учёные-кондитеры,
они же --- задачи с сайта informatics. Внедрённый рецепт приводит в восторг
инвесторов (\emph{«Это же конфета с нейроинтерфейсом! Берём!»}),
и они засыпают фабрику деньгами и ускоряют производство.

Правда, пока конструкторское бюро возится с рецептом, оно нагло ворует
электричество у конвейера --- производство замедляется. А за отдельную плату
штатный промышленный шпион готов вынести с чужой фабрики
\emph{инсайдерскую информацию} (тест к задаче). Шпион берёт предоплату,
работает без гарантий и чеков не выдаёт.

К дедлайну, объявленному Всемирной Кондитерской Комиссией, побеждает
фабрика с самым большим складом конфет. Остальных ждёт слияние,
поглощение и переработка в повидло.

\section*{Как всё устроено}

\begin{enumerate}[leftmargin=1.6em,itemsep=2pt]

\item \textbf{Стартовый капитал.} Каждая фабрика начинает с одинаковыми
\term{Запасами} конфет и одинаковой \term{Скоростью} производства.
Конвейер не останавливается ни на секунду: каждую секунду игрового времени
склад пополняется на величину текущей \term{Скорости}. Даже если вся
команда ушла пить чай.

\item \textbf{Каталог рецептов.} Всем фабрикам доступен один и тот же
набор задач: несколько уровней сложности, в каждом уровне --- несколько
задач. Но порядок задач внутри уровня у каждой фабрики свой, секретный
и перемешанный: подглядывать в табло конкурента бессмысленно ---
их ячейка № 5 и ваша ячейка № 5 --- почти наверняка разные задачи.
Промышленный шпионаж через плечо не работает, мы проверяли.

\item \textbf{Покупка рецепта (задачи).} Условие задачи засекречено,
пока рецепт не куплен: платите \term{Цену задачи} со склада --- получаете
ссылку на условие. Купить задачу команда может \emph{сама}, кнопкой
на своей странице. Касса сдачи не даёт и возвратов не принимает:
потраченные на покупку конфеты не возвращаются \emph{никогда},
даже по заявлению в трёх экземплярах.

\item \textbf{Плата за прогресс.} С момента покупки и до сдачи рецепта
конструкторское бюро ворует мощности: \term{Скорость} снижается
на \term{Нагрузку от задачи}. Хотите купить всё и сразу --- считайте:
если конфет на складе меньше цены, а конвейер и так еле дышит,
фабрика сделку не одобрит.

\item \textbf{Инсайдерская информация (тесты).} Любое количество,
по \term{Цене подсказки} за штуку, приобретается через ведущего ---
шпион работает только с директором лично. Деньги за инсайд
не возвращаются, качество инсайда обсуждению не подлежит (пункт 1
трудового договора шпиона).

\item \textbf{Внедрение рецепта (решение задачи).} Решение сдаётся
на informatics. Как только вердикт --- \emph{OK} или \emph{Зачтено},
фабрика автоматически получает всё и сразу:
\begin{itemize}[itemsep=0pt]
	\item склад пополняется на \term{Бонус запасы} --- инвесторы ликуют;
	\item конструкторское бюро возвращает украденные мощности
	(\term{Нагрузка} снова прибавляется к скорости);
	\item и сверху --- \term{Бонус скорость}: конвейер разгоняется навсегда.
\end{itemize}
Моментом внедрения считается \emph{время отправки посылки}, а не время
её проверки. Если тестирующая система задумалась --- это её трудности,
а не ваши.

\item \textbf{Посылка на флажке.} Решение, отправленное до дедлайна,
но проверенное после, --- засчитывается. Отправленное после дедлайна ---
уходит в архив Всемирной Кондитерской Комиссии с пометкой
\emph{«надо было раньше»}.

\item \textbf{Без самодеятельности.} Сдавать некупленную задачу
бесполезно: рецепт, украденный без оплаты, автоматика не внедряет,
а отправляет ведущему на разбирательство. Ведущий, конечно, может
и зачесть --- но сперва очень выразительно на вас посмотрит.

\item \textbf{Победа.} В момент дедлайна побеждает фабрика с наибольшими
\term{Запасами}. Скорость, амбиции и красивые глаза в зачёт не идут ---
только конфеты на складе.

\item \textbf{Ведущий --- окончательный авторитет.} Он видит все сделки
в журнале, может продлить игру, исправить ошибочную запись и вручную
зачесть решение, если informatics прилёг отдохнуть. Спорить с ведущим
можно, но статистика споров не на вашей стороне: 0 побед из 0 возможных
у предыдущих составов.

\end{enumerate}

\section*{Параметры уровней}

Цены и бонусы зависят от уровня сложности задачи и объявляются
на табло игры (таблица \emph{«Параметры игры»}). Чем сложнее рецепт,
тем дороже инсайд и тем щедрее инвесторы.

\begin{center}
\begin{tabular}{lccc}
\toprule
 & \textbf{Уровень 1} & \textbf{Уровень 2} & \textbf{Уровень 3} \\
\midrule
Цена задачи        & 12\,000 & 12\,000 & 12\,000 \\
Цена подсказки     & 3\,000  & 7\,000  & 10\,000 \\
Нагрузка от задачи & 2       & 2       & 2       \\
Бонус запасы       & 12\,000 & 25\,000 & 50\,000 \\
Бонус скорость     & 4       & 7       & 11      \\
\bottomrule
\end{tabular}

\small\emph{(значения по умолчанию; точные параметры вашей игры --- на табло)}
\end{center}

\section*{Шпаргалка по цветам табло}

\begin{center}
\begin{tabular}{ll}
\toprule
\colorbox{gray!30}{\strut ~серая~}   & задача не куплена: рецепт в сейфе, условие скрыто \\
\colorbox{yellow!60}{\strut ~жёлтая~} & куплена: бюро работает, конвейер страдает \\
\colorbox{green!40}{\strut ~зелёная~} & решена: инвесторы аплодируют стоя \\
\bottomrule
\end{tabular}
\end{center}

\vspace{1em}
\begin{center}
\emph{Удачного майнинга конфет! И помните: возвратов нет.}
\end{center}

\end{document}
